\chapter[Proposta]{Proposta}
\label{cap:proposta}

Este capítulo apresenta o artefato proposto. A Seção~\ref{sec:requisitos}
define os requisitos da arquitetura. A Seção~\ref{sec:arquitetura} descreve a
arquitetura por meio das camadas do \textit{Data Lakehouse} e dos papéis que
cada componente desempenha, sem se prender a ferramentas específicas. A
Seção~\ref{sec:caso} apresenta a instanciação no caso de prevenção de
paralisação de obras do MCMV, incluindo as ferramentas escolhidas para cada
papel.

% ---------------------------------------------------------------------
\section{Requisitos da arquitetura}
\label{sec:requisitos}

Os requisitos derivam de duas fontes: os riscos e princípios levantados no
Capítulo~\ref{cap:referencial} e as regras de engenharia que a plataforma
Gov Hub Cidades já adota para todos os seus dados. Essas regras são:

\begin{enumerate}
	\item \textbf{O cálculo mora na camada de transformação.} O dado chega pelo
	orquestrador, é transformado entre as camadas do \textit{lakehouse}, e a
	camada de visualização apenas filtra;
	\item \textbf{Não se cria camada nova para acomodar complexidade.} Antes de
	criar um conceito, verifica-se se ele precisa existir;
	\item \textbf{O \textit{lake} é a interface de arquivos.} Todo dado entra
	por uma zona de pouso de arquivos, que espelha a fonte, e só dali é lido
	pela camada bronze.
\end{enumerate}

A Tabela~\ref{tab:requisitos} lista os requisitos resultantes.

\begin{table}[htb]
	\centering
	\caption{Requisitos da arquitetura de MLOps.}
	\label{tab:requisitos}
	\small
	\begin{tabularx}{\textwidth}{p{0.8cm}X>{\raggedright\arraybackslash}p{4.2cm}}
		\toprule
		\textbf{ID} & \textbf{Requisito} & \textbf{Origem} \\
		\midrule
		R1 & Toda feature é calculada uma única vez, por uma transformação
		declarativa e versionada, e usada tanto no treino quanto no escore. &
		\citeonline{sculley2015}; \citeonline{polyzotis2018}; regra 1 \\
		R2 & Toda feature de uma observação no mês $t$ usa apenas dados
		disponíveis até $t$, e isso é verificado por teste automatizado. & \citeonline{kaufman2012} \\
		R3 & Todo resultado de treino pode ser refeito a partir da versão do
		código e da versão do dataset. & \citeonline{kreuzberger2023} (P3, P4); \citeonline{shankar2024} \\
		R4 & Todo experimento registra parâmetros, métricas, artefatos e a
		origem dos dados. & \citeonline{zaharia2018}; \citeonline{kreuzberger2023} (P7) \\
		R5 & Um modelo só substitui o vigente se superar o vigente e a regra de
		referência, na mesma métrica e no mesmo período. & \citeonline{studer2021}; \citeonline{breck2017} \\
		R6 & Distribuição das features e desempenho realizado são monitorados
		continuamente. & \citeonline{gama2014}; \citeonline{kreuzberger2023} (P8) \\
		R7 & Tabelas de ML passam pelas mesmas camadas, testes e catálogo que
		qualquer outra tabela da plataforma. & \citeonline{ferreira2025}; regras 1 e 2 \\
		R8 & Adicionar um segundo modelo não exige alterar a estrutura do
		primeiro. & \citeonline{amershi2019} \\
		\bottomrule
	\end{tabularx}
	\fonte{Elaborado pelos autores.}
\end{table}

% ---------------------------------------------------------------------
\section{Arquitetura proposta}
\label{sec:arquitetura}

\subsection{Visão geral}

A arquitetura opera sobre um \textit{Data Lakehouse} \cite{armbrust2021}
organizado nas três camadas da arquitetura medalhão
(Seção~\ref{sec:plataformas}). A ideia central é estender o
\textit{lakehouse} existente em vez de construir uma plataforma de ML
paralela: cada artefato do ciclo de vida de um modelo (painel, rótulo,
features, escores e métricas de monitoramento) ocupa uma camada definida e
passa pelos mesmos testes e pelo mesmo catálogo das demais tabelas.

A arquitetura é descrita por \textbf{papéis}, e não por ferramentas. Cada papel
define uma responsabilidade que pode ser cumprida por diferentes tecnologias.
Essa separação permite reaproveitar a arquitetura em plataformas com pilhas
tecnológicas diferentes. As ferramentas que materializam cada papel no caso
estudado estão na Seção~\ref{sec:ferramentas-caso}.

A Figura~\ref{fig:arquitetura} apresenta a visão geral.

\begin{figure}[htb]
	\centering
	\resizebox{\textwidth}{!}{%
	\begin{tikzpicture}[
		box/.style={draw, rounded corners, align=center, minimum height=1.3cm,
			text width=3.2cm, font=\footnotesize},
		bronze/.style={box, fill=orange!18},
		prata/.style={box, fill=gray!15},
		ouro/.style={box, fill=yellow!28},
		ml/.style={box, fill=blue!8},
		arr/.style={-{Stealth[length=2mm]}, thick},
		elbl/.style={font=\scriptsize\itshape, fill=white, inner sep=1pt},
		zone/.style={draw, dashed, rounded corners, inner sep=0.3cm},
		lbl/.style={font=\scriptsize\itshape, anchor=south west}
	]
		% Entrada
		\node[box] (ingest) at (0,5.2) {Orquestrador\\ingestão das fontes};

		% Camadas medalhão
		\node[bronze] (bronze) at (0,2.6)  {\textbf{Bronze}\\dado bruto da fonte};
		\node[prata]  (prata)  at (0,0)    {\textbf{Prata}\\dado limpo; painel e rótulo};
		\node[ouro]   (ouro)   at (0,-2.6) {\textbf{Ouro}\\features, escores e monitoramento};

		% Consumo e governança
		\node[box] (vis) at (0,-5.2) {Camada de\\visualização};
		\node[box] (cat) at (-4.6,0) {Catálogo de dados};

		% Extensão de MLOps
		\node[ml] (gate)   at (4.8,2.6)  {Orquestrador\\promoção com \textit{gate}};
		\node[ml] (reg)    at (9.4,2.6)  {Registro de experimentos e modelos};
		\node[ml] (escore) at (4.8,0)    {Ambiente de modelagem\\escore mensal};
		\node[ml] (treino) at (9.4,0)    {Ambiente de modelagem\\treino e avaliação};
		\node[ml] (export) at (4.8,-2.6) {Orquestrador\\export do dataset};
		\node[ml] (ds)     at (9.4,-2.6) {Armazenamento de objetos\\datasets versionados};

		% Fluxo de dados
		\draw[arr] (ingest) -- (bronze);
		\draw[arr] (bronze) -- node[elbl, left] {transformação} (prata);
		\draw[arr] (prata) -- node[elbl, left] {transformação} (ouro);
		\draw[arr] (ouro) -- (vis);
		\draw[arr, dashed] (cat) -- (prata);

		% Ciclo de ML
		\draw[arr] (ouro) -- (export);
		\draw[arr] (export) -- (ds);
		\draw[arr] (ds) -- (treino);
		\draw[arr] (treino) -- (reg);
		\draw[arr] (reg) -- (gate);
		\draw[arr] (gate) -- (escore);
		\draw[arr] (ouro.north east) -- node[elbl, sloped, below] {features} (escore.south west);
		\draw[arr] (escore.north west) -- node[elbl, sloped, above] {escores} (bronze.south east);

		\begin{scope}[on background layer]
			\node[zone, fit=(bronze)(prata)(ouro)] (z1) {};
			\node[zone, fit=(export)(ds)(treino)(reg)(gate)(escore), draw=blue!60] (z2) {};
		\end{scope}
		\node[lbl, anchor=north east, align=right] at (z1.north west) {\textit{Data Lakehouse}\\(camadas medalhão)};
		\node[lbl, text=blue!60!black] at (z2.north west) {extensão de MLOps};
	\end{tikzpicture}%
	}
	\caption{Visão geral da arquitetura de MLOps proposta.}
	\label{fig:arquitetura}
	\fonte{Elaborado pelos autores.}
\end{figure}

\subsection{Camadas do \textit{lakehouse} no ciclo de ML}

A Tabela~\ref{tab:camadas} mostra o que cada camada da arquitetura medalhão
guarda, em geral e no ciclo de ML.

\begin{table}[htb]
	\centering
	\caption{Conteúdo de cada camada do \textit{lakehouse}.}
	\label{tab:camadas}
	\small
	\begin{tabularx}{\textwidth}{>{\raggedright\arraybackslash}p{1.6cm}XX}
		\toprule
		\textbf{Camada} & \textbf{Conteúdo geral} & \textbf{Conteúdo de ML} \\
		\midrule
		Bronze & Cópia da zona de pouso, ainda bruta e sem regra de negócio.
		Ponto de reprocessamento das demais camadas & Escores, lidos dos
		arquivos que o ambiente de modelagem grava na zona de pouso, um por
		competência \\
		Prata & Dado limpo, tipado e harmonizado entre fontes & Painel
		unificado (uma linha por unidade e mês) e rótulo \\
		Ouro & Tabelas modeladas para consumo & Features calculadas
		\textit{as of} $t$, escores consolidados e métricas de \textit{drift}
		e de desempenho realizado \\
		\bottomrule
	\end{tabularx}
	\fonte{Elaborado pelos autores.}
\end{table}

Antes da camada bronze há uma zona de pouso de arquivos, que guarda o dado no
formato da fonte e serve de interface entre a ingestão e o
\textit{lakehouse} (regra 3). Ela não é uma camada medalhão: nenhuma
transformação de negócio acontece ali.

Dois pontos merecem destaque. Primeiro, os escores voltam ao
\textit{lakehouse} como qualquer outra fonte: são gravados como arquivos na
zona de pouso, lidos pela camada bronze e chegam à camada ouro pelas mesmas
transformações e testes. O ambiente de modelagem
nunca escreve diretamente na camada ouro. Segundo, os datasets congelados de
cada treino e os artefatos de modelo ficam em uma área própria do
armazenamento de objetos, fora das camadas medalhão. Eles são saídas do
treino, e não etapas de transformação de dados, por isso não justificam uma
camada nova (regra 2).

\subsection{Papéis e responsabilidades}

A principal decisão da arquitetura é qual papel assume cada responsabilidade
do ciclo de vida de ML (QE1). A Tabela~\ref{tab:responsabilidades} resume essa
divisão. A coluna ``Não faz'' é tão importante quanto a de responsabilidades,
porque é ela que impede a duplicação de lógica entre componentes.

\begin{table}[htb]
	\centering
	\caption{Papéis da arquitetura e suas responsabilidades.}
	\label{tab:responsabilidades}
	\small
	\begin{tabularx}{\textwidth}{>{\raggedright\arraybackslash}p{3.2cm}X>{\raggedright\arraybackslash}p{3.0cm}}
		\toprule
		\textbf{Papel} & \textbf{Responsabilidade} & \textbf{Não faz} \\
		\midrule
		Orquestrador & Agendar e encadear ingestão, export de datasets, treino,
		promoção e escore; aplicar o \textit{gate} de promoção & Conter cálculo
		sobre dados \\
		Camada de transformação & Calcular painel, rótulo, features
		\textit{as of} $t$, leitura dos escores e métricas de monitoramento;
		executar testes de dados e de vazamento & Treinar modelos \\
		Motor analítico & Materializar e consultar as camadas bronze, prata e
		ouro & --- \\
		Armazenamento de objetos & Guardar a zona de pouso, os datasets
		congelados e os artefatos de modelo & --- \\
		Ambiente de modelagem & Treinar, avaliar e escorar, lendo apenas
		tabelas da camada ouro & Recalcular features \\
		Registro de experimentos e modelos & Registrar parâmetros, métricas,
		artefatos, versões de modelo e o estágio de cada versão & Decidir a
		promoção \\
		Catálogo de dados & Registrar tabelas, responsáveis e linhagem de todas
		as camadas & --- \\
		Camada de visualização & Exibir o ranking de risco e os indicadores de
		monitoramento & Calcular métricas \\
		\bottomrule
	\end{tabularx}
	\fonte{Elaborado pelos autores.}
\end{table}

Em \textit{lakehouses} que usam formatos de tabela abertos sobre armazenamento
de objetos, o motor analítico e o armazenamento de objetos podem coincidir
\cite{armbrust2021}. A arquitetura os separa porque nem toda plataforma
governamental adota esse formato.

Como a mesma transformação gera o dataset de treino e o de escore, não há duas
implementações da mesma feature (R1), o que elimina por construção a principal
causa de desalinhamento entre treino e produção \cite{polyzotis2018}. Uma
feature nova é sempre uma transformação versionada da camada ouro, testada e
catalogada como qualquer outra tabela (R7).

\subsection{Ciclo de vida}

Cada modelo percorre as fases do CRISP-DM \cite{wirth2000}, e cada fase tem
um lugar definido na arquitetura (Tabela~\ref{tab:ciclo}). A análise
exploratória vem antes de qualquer transformação de ML: nenhuma tabela nova é
criada antes de a exploração mostrar que ela é necessária, em respeito à
regra 2 da plataforma.

\begin{table}[htb]
	\centering
	\caption{Fases do CRISP-DM e sua localização na arquitetura.}
	\label{tab:ciclo}
	\small
	\begin{tabularx}{\textwidth}{>{\raggedright\arraybackslash}p{3.0cm}>{\raggedright\arraybackslash}p{4.8cm}X}
		\toprule
		\textbf{Fase} & \textbf{Onde mora} & \textbf{Saída} \\
		\midrule
		1. Entendimento do problema & Documentação do modelo, versionada com
		o código & Definição do evento, usuário, antecedência, capacidade e
		referência de comparação \\
		2. Entendimento dos dados & Análise exploratória, que lê a camada prata
		e nunca vai para produção & Relatório com decisões de rótulo, janela e
		features \\
		3. Preparação & Camada de transformação (prata e ouro) & Painel, rótulo
		e features \textit{as of} $t$, com teste de vazamento \\
		4. Modelagem & Ambiente de modelagem e registro de experimentos &
		Experimentos registrados \\
		5. Avaliação & Ambiente de modelagem & Métricas contra a regra de
		referência, em corte temporal \\
		6. Implantação e operação & Orquestrador, camada de transformação e
		camada de visualização & Escore mensal, retreino com \textit{gate},
		monitoramento \\
		\bottomrule
	\end{tabularx}
	\fonte{Elaborado pelos autores.}
\end{table}

\subsection{Reprodutibilidade}

O treino nunca lê uma base extraída manualmente. Ele lê a camada ouro e
congela uma cópia versionada no armazenamento de objetos, acompanhada de um
manifesto com a versão do código (\textit{commit}), a versão das
transformações e as competências incluídas. A versão do dataset é registrada
no experimento. Assim, qualquer resultado pode ser refeito a partir da versão
do código e da versão do dataset (R3, QE2), o que será verificado pelo
critério AV3.

\subsection{Promoção e operação}

A promoção de um modelo é controlada por um \textit{gate} aplicado pelo
orquestrador: uma nova versão só passa a gerar o escore oficial se, no período
de validação temporal, superar tanto a versão vigente quanto a regra de
referência no \textit{recall@k} (R5). O escore mensal é gravado na camada
bronze e promovido até a camada ouro pela camada de transformação, o que
mantém o ranking sob os mesmos testes e o mesmo catálogo das demais tabelas.

O monitoramento também fica na camada de transformação (R6): tabelas da
camada ouro comparam a distribuição das features de cada competência com a do
dataset de treino e, à medida que o horizonte de predição se cumpre, calculam
o desempenho realizado dos escores passados. A camada de visualização exibe
esses indicadores.

\pendente{Definir os limiares de alerta de \textit{drift} (por exemplo, PSI
ou teste de Kolmogorov--Smirnov por feature) e a política de retreino
(periódico ou disparado por alerta).}

% ---------------------------------------------------------------------
\section{Instanciação: prevenção de paralisação de obras do MCMV}
\label{sec:caso}

\subsection{Problema}

Para cada obra \textbf{em execução} no mês $t$, estima-se a probabilidade de
ela estar \textbf{paralisada} em algum mês até $t+h$. A saída é um ranking
mensal de risco, para que a fiscalização priorize visitas antes de a obra
parar.

A unidade de análise é a obra (\texttt{id\_empreendimento}, que acompanha o
empreendimento por todas as fases) observada em um mês de competência. O
escopo inclui as modalidades com execução de obra a acompanhar: Fundo de
Arrendamento Residencial (FAR), Fundo de Desenvolvimento Social (FDS), Rural,
Pró-Moradia e Financiado (FGTS). Modalidades que são contratos ou propostas,
sem execução de obra, ficam de fora.

\subsection{Dados disponíveis}

O painel histórico do FAR, do FDS e do Rural está em tabelas prata de
histórico por empreendimento, uma por modalidade, com grão de obra por mês,
de dezembro de 2019 a julho de 2026. São 5.555 obras no FAR, 1.021 no FDS
e 11.122 no Rural. A situação da obra é harmonizada em uma coluna canônica com
os valores \texttt{em\_obras}, \texttt{paralisada}, \texttt{concluida},
\texttt{cancelada} e \texttt{nao\_iniciada}. O FGTS tem tabelas próprias, com
percentual previsto e realizado, data da última vistoria e motivos de
paralisação.

A fonte da situação muda duas vezes ao longo da série: SFTP da Caixa até
novembro de 2024, sistema da SNH de junho de 2024 a março de 2026, e planilha
de acompanhamento mensal a partir de abril de 2026. No período de sobreposição,
as duas primeiras fontes concordam em 96 a 98\% dos casos. Ainda assim, um
terço das paralisações do FAR ocorre exatamente na troca de fonte, e dentro do
período SNH quase não surgem paralisações novas. No período SFTP há cerca de
1.500 transições de \texttt{em\_obras} para \texttt{paralisada}, a maioria no
Rural.

Essas observações indicam que o rótulo só é confiável no período SFTP, e que
a troca de fonte é uma mudança de distribuição conhecida
(Seção~\ref{sec:riscos}) que a preparação dos dados precisa tratar
explicitamente.

\subsection{Referência de comparação}

A regra vigente (\texttt{ouro\_reloginho\_indicadores\_gargalo\_desempenho})
sinaliza uma obra quando ela passa 90 dias sem atualização, apresenta
diferença superior a 10 pontos percentuais entre execução física e financeira
ou tem menos de 30\% de execução. O modelo só se justifica se, com o mesmo
número de visitas, antecipar mais paralisações ou antecipá-las mais cedo.

\subsection{Métrica e modelagem}

A métrica principal é o \textit{recall@k}, em que $k$ é a capacidade mensal de
fiscalização. Como métricas secundárias, usam-se \textit{precision@k} e a área
sob a curva precisão-revocação, mais informativa que a curva ROC em classes
desbalanceadas \cite{saito2015}. A divisão entre treino, validação e teste é
temporal.

Como primeira família de modelos, propõe-se \textit{gradient boosting} com
LightGBM \cite{ke2017}, com explicações individuais por valores SHAP
\cite{lundberg2017}, para que a fiscalização saiba por que uma obra foi
priorizada. Uma formulação alternativa por análise de sobrevivência
\cite{wang2019survival} será considerada se a EDA mostrar que o tempo até a
paralisação é mais informativo que o evento binário. O modelo será documentado
em um \textit{model card} \cite{mitchell2019}.

\subsection{Materialização da arquitetura}
\label{sec:ferramentas-caso}

A Tabela~\ref{tab:ferramentas-caso} mostra a ferramenta que cumpre cada papel
da arquitetura no Gov Hub Cidades e onde ela fica no repositório. A maior
parte das ferramentas já fazia parte da plataforma. As novas são o MLflow e as
bibliotecas de modelagem.

\begin{table}[htb]
	\centering
	\caption{Ferramentas que materializam cada papel no Gov Hub Cidades.}
	\label{tab:ferramentas-caso}
	\small
	\begin{tabularx}{\textwidth}{>{\raggedright\arraybackslash}p{3.4cm}>{\raggedright\arraybackslash}X>{\raggedright\arraybackslash}p{4.8cm}}
		\toprule
		\textbf{Papel} & \textbf{Ferramenta} & \textbf{Local no repositório} \\
		\midrule
		Orquestrador & Apache Airflow \cite{airflow} & \texttt{dags/ml/} \\
		Camada de transformação & dbt \cite{dbt} & \texttt{dbt/mcid/models/} \texttt{ml\_<modelo>\_dbt/} \\
		Motor analítico & PostgreSQL com a extensão pg\_duckdb & --- \\
		Armazenamento de objetos & RustFS, compatível com S3 & prefixos \texttt{raw/}, \texttt{staging/} e \texttt{ml/} \\
		Ambiente de modelagem & Python, com pandas, scikit-learn, LightGBM e SHAP & \texttt{ml/<modelo>/} \\
		Registro de experimentos e modelos & MLflow \cite{zaharia2018} & artefatos em \texttt{ml/mlflow/} \\
		Catálogo de dados & OpenMetadata \cite{openmetadata} & --- \\
		Camada de visualização & Apache Superset & --- \\
		Análise exploratória & Jupyter & \texttt{ml/<modelo>/notebooks/} \\
		Versionamento e integração contínua & Git e GitHub Actions & --- \\
		\bottomrule
	\end{tabularx}
	\fonte{Elaborado pelos autores.}
\end{table}

Na plataforma (Figura~\ref{fig:govhub}), a zona de pouso são os prefixos
\texttt{raw/} e \texttt{staging/} do armazenamento de objetos, e as
camadas bronze, prata e ouro são tabelas materializadas pelo dbt no
PostgreSQL, que lê o Parquet da \texttt{staging/} pela extensão pg\_duckdb.
Os escores do modelo seguem o mesmo caminho, a partir de um prefixo próprio
do armazenamento de objetos (\texttt{ml/escores/}).

\subsection{Decisões em aberto}

As seguintes decisões dependem de alinhamento com o MCID e precisam ser
tomadas antes da fase de preparação dos dados:

\begin{itemize}
	\item \textbf{Horizonte $h$}: com quantos meses de antecedência a
	fiscalização consegue agir. Proposta: seis meses, com análise de
	sensibilidade para três e doze;
	\item \textbf{Capacidade $k$}: quantas obras a fiscalização visita por mês;
	\item \textbf{Retomada}: se uma obra que retoma e paralisa de novo conta
	como um segundo evento;
	\item \textbf{Usuário do alerta}: área técnica do MCID, agente financeiro
	ou ambos, o que define o que o escore precisa explicar.
\end{itemize}

\pendente{Registrar aqui as decisões tomadas sobre $h$, $k$, retomada e
usuário do alerta, com data e responsável.}
